%%%%%%%%%%%%%%%%%%%%%%%%%%%%%%%%%%%%%%%%%%%%%%%%%%%%%%%%%%%%%%%%%%%%%%%%%%%%%%%%

\begin{english}
\thispagestyle{empty}
\vspace*{25pt}
\centerline{\bfseries \abstractname}
Weather forecasting is essential for many societal and economic activities, including agriculture, transportation, disaster prevention, and energy management. Traditional forecasting systems rely primarily on Numerical Weather Prediction (NWP), which uses physical models to simulate atmospheric dynamics. Although these models provide highly accurate predictions, they require considerable computational resources and specialized infrastructure.
Recent advances in artificial intelligence and machine learning have opened new possibilities for data-driven weather forecasting. Deep learning models are capable of learning complex spatial and temporal relationships directly from large meteorological datasets.
This project investigates the application of deep learning methods for short-term weather forecasting using the ERA5 reanalysis dataset. Several neural network architectures will be explored, including convolutional neural networks and encoder–decoder models such as U-Net. The models are trained using PyTorch and evaluated using quantitative metrics and visualization techniques.
The objective of this work is to assess the ability of deep learning models to capture atmospheric patterns and generate reliable weather predictions.
\vspace*{25pt}

\begin{description}
\item [Keywords:] \mbox{} \\[-20pt]
\begin{itemize}
  \item Deep Learning
  \item Numerical Weather Prediction
  \item ERA5 Reanalysis
  \item Convolutional Neural Networks
  \item Encoder–Decoder Architecture
  \item Temperature Forecasting
  \item Spatio-Temporal Modeling

\end{itemize}
\end{description}
\end{english}

%%%%%%%%%%%%%%%%%%%%%%%%%%%%%%%%%%%%%%%%%%%%%%%%%%%%%%%%%%%%%%%%%%%%%%%%%%%%%%%%
